% =====================================================================
%  Wspólna preambuła szablonów pism ZUS (świadczenia rodzicielskie 2026)
%  Silnik: XeLaTeX albo Tectonic (pdflatex NIE obsłuży polskich znaków w fontspec).
%  Kompilacja:  tectonic -X compile plik.tex
% =====================================================================
\documentclass[11pt,a4paper]{article}

\usepackage{fontspec}
\usepackage{polyglossia}
\setmainlanguage{polish}
\setotherlanguage{english}
% font domyślny (Latin Modern) — Tectonic go dostarcza

\usepackage[margin=2.5cm,a4paper]{geometry}
\usepackage{graphicx}
\usepackage{xcolor}
\usepackage{array}
\usepackage{enumitem}
\usepackage{fancyhdr}
\usepackage{hyperref}

\definecolor{zus}{HTML}{007A33}      % zieleń ZUS
\definecolor{szary}{HTML}{6B7280}
\definecolor{ramka}{HTML}{9CA3AF}

\hypersetup{colorlinks=true, linkcolor=zus, urlcolor=zus, pdftitle={Szablony pism ZUS}}

% --- pole do wypełnienia odręcznego -----------------------------------
% \Pole{nazwa} rysuje kropkowaną linię z podpisem pola pod spodem.
% Pole do wypełnienia: kropkowana linia, pod nią drobny podpis pola.
\newcommand{\Pole}[1]{%
  \textcolor{szary}{\dotfill}\,%
  {\scriptsize\textcolor{szary}{\textsc{#1}}}%
}

% --- blok warunkowy: treść pokazywana tylko, gdy pole prawdziwe --------
\newenvironment{warunek}[1]{%
  \par\smallskip\noindent\begin{minipage}{\linewidth}%
  \color{black}%
}{%
  \end{minipage}\par\smallskip%
}

% --- nagłówek pisma ---------------------------------------------------
\newcommand{\NaglowekPisma}[1]{%
  \par\vspace{0.8em}%
  \begin{center}%
    {\large\bfseries\MakeUppercase{#1}}%
  \end{center}%
  \vspace{0.4em}%
}

% --- blok nadawcy (pierwsze linie pisma) -------------------------------
\newenvironment{nadawca}{%
  \begin{flushleft}\small
}{%
  \end{flushleft}\normalsize
}

% --- stopka z numerem strony -----------------------------------------
\pagestyle{fancy}
\fancyhf{}
\renewcommand{\headrulewidth}{0pt}
\fancyfoot[C]{\textcolor{szary}{\small\thepage}}

% --- metadane szablonu na początku pisma ------------------------------
\newcommand{\Metadane}[3]{%
  \begin{flushright}\footnotesize\textcolor{szary}%
    \textbf{#1}\\[2pt] {\itshape #2}\\[2pt] #3%
  \end{flushright}\vspace{0.5em}%
}

\setlength{\parindent}{0pt}
\setlength{\parskip}{0.6em}

\begin{document}

\begin{flushright}
{\footnotesize\color{szary}
\textbf{Szablon P03}\\
{\itshape Wniosek o zasiłek macierzyński za okres urlopu rodzicielskiego (formularz ZUR)}\\
art. 182\textsuperscript{1}d k.p.; art. 29--31 ustawy zasiłkowej -- Formularz: ZUR (do 24.10.2026), potem ZAM}\
\end{flushright}

\vspace{0.5em}
{\color{ramka}\hrule height 0.4pt}
\vspace{0.8em}

[FORMULARZ ZUS --- ZUR / ZAM]\\[-2pt]
Wniosek o zasiłek macierzyński za okres urlopu rodzicielskiego

Wypełnij ten wniosek, jeśli ubiegasz się o wypłatę zasiłku macierzyńskiego
za okres urlopu rodzicielskiego albo urlopu na warunkach urlopu macierzyńskiego.

ZAKRES WNIOSKU
Proszę o wypłatę zasiłku macierzyńskiego za okres od \Pole{urlop od} do \Pole{urlop do}.

Wnioskuję o zasiłek macierzyński za okres urlopu rodzicielskiego z powodu:
[ ] urodzenia się dziecka/dzieci
[ ] przyjęcia na wychowanie dziecka/dzieci w celu przysposobienia
[ ] przyjęcia na wychowanie dziecka/dzieci w ramach rodziny zastępczej
[ ] przerwania przez matkę dziecka okresu wypłaty zasiłku macierzyńskiego
[ ] posiadania przez matkę orzeczenia o niezdolności do samodzielnej egzystencji
[ ] porzucenia dziecka przez matkę
[ ] śmierci matki dziecka/dzieci
[ ] podjęcia przez matkę pracy w wymiarze nie niższym niż połowa pełnego etatu
[ ] urodzenia dziecka/dzieci po \Pole{data porodu}

DANE DZIECKA/DZIECI
\Pole{dziecko imie}, \Pole{data porodu}, \Pole{adres miejscowosc}

DANE DRUGIEGO RODZICA
\Pole{drugi rodzic}, PESEL \Pole{pesel}

ŚWIADCZENIA PROSZĘ PRZEKAZYWAĆ
[ ] w Polsce na rachunek bankowy
[ ] za granicę --- na rachunek bankowy (oprócz USA)

\vfill
{\footnotesize\textcolor{szary}\itshape Zgodnie z art. 30 ustawy zasiłkowej wniosek o zasiłek za okres urlopu rodzicielskiego składa się w terminie do 21 dni po porodzie, jeśli chce się zasiłek za cały okres; późniejsze wnioski --- za okres od dnia złożenia. }

\end{document}
